% !TeX spellcheck = en_US
\documentclass[12pt,a4paper]{amsart}

\usepackage[T1]{fontenc}
\usepackage{lmodern}
\usepackage{a4wide,fullpage,setspace}
\usepackage{amsmath,amssymb,mathtools}
\usepackage[shortlabels]{enumitem}
\usepackage[numbers,sort&compress]{natbib}
\usepackage{tikz-cd}
\definecolor{dark-red}{rgb}{0.5,0.15,0.15}
\usepackage[colorlinks=true,linkcolor=black,citecolor=dark-red,urlcolor=dark-red]{hyperref}
\setstretch{1.1}
\raggedbottom
%\flushbottom
\swapnumbers

%\usepackage{lineno}
%\linenumbers

%%%%%%%%%%%%%%%%%%%%%%%%%%%%%%%%%%%
%use of renewcommand too dangerous
\makeatletter
\let\P\@undefined
\let\leq\@undefined
\let\geq\@undefined
\let\vec\@undefined
\let\phi\@undefined
\let\epsilon\@undefined
\makeatother
\newcommand{\leq}{\leqslant}
\newcommand{\geq}{\geqslant}
\newcommand{\vec}{\overrightarrow}
\newcommand{\phi}{\varphi}
\newcommand{\epsilon}{\varepsilon}
%%%%%%%%%%%%%%%%%%%%%%%%%%%%%%%%%%%

\newcommand{\Top}{\mathop{\mathsf{Top}}}
\newcommand{\TOP}{\mathop{\mathsf{TOP}}}
\newcommand{\Pcat}{\mathcal P}
\newcommand{\G}{\mathcal G}
\newcommand{\M}{\mathcal M}
\newcommand{\One}{\mathbf 1}
\newcommand{\PSp}{[\Pcat^{\mathrm{op}},\Top]_0}
\newcommand{\PFlow}{\Pcat\mathop{\mathsf{Flow}}}
\newcommand{\Path}{\mathbb P}
\newcommand{\Ffree}{\mathbb F}
\newcommand{\id}{\operatorname{id}}
\newcommand{\Obj}{\operatorname{Obj}}
\newcommand{\Pmult}{\Pcat_{\times}}
\newcommand{\cocartesian}{\arrow[lu, phantom, "\ulcorner"{font=\Large}, pos=0]}
\newcommand{\cartesian}{\arrow[phantom,rd,"\lrcorner"{font=\Large},pos=0]}
%\newcommand{\liminj}{\mathop{\mathrm{colim}}}
\newcommand{\liminj}{\varinjlim}
\newcommand{\limproj}{\varprojlim}
%\makeatletter
%%
%\def\varholim@#1#2{%
%	\vtop{\m@th\ialign{##\cr
%			\hfil\(#1\operator@font holim\)\hfil\cr
%			\noalign{\nointerlineskip\kern1.5\ex@}#2\cr
%			\noalign{\nointerlineskip\kern-\ex@}\cr}}%
%}
%\def\holimproj{%
%	\mathop{\mathpalette\varholim@{\leftarrowfill@\textstyle}}\nmlimits@
%}
%\def\holiminj{%
%	\mathop{\mathpalette\varholim@{\rightarrowfill@\textstyle}}\nmlimits@
%}
%\makeatother

%\newtheorem*{thmN}{Theorem}
%\newtheorem*{quN}{Question}
\newtheorem{thm}{Theorem}[section]
\newtheorem{lem}[thm]{Lemma}
\newtheorem{proposition}[thm]{Proposition}
\newtheorem{cor}[thm]{Corollary}
\theoremstyle{definition}
\newtheorem{definition}[thm]{Definition}
\newtheorem{remark}[thm]{Remark}


\title{Left properness of Moore flows}
\author[P. Gaucher]{Philippe Gaucher}
\address{Universit\'e Paris Cit\'e, CNRS, IRIF, F-75013, Paris, France}
\urladdr{\url{https://www.irif.fr/~gaucher}}
\subjclass[2020]{Primary 18N40; Secondary 55U35, 55U40, 18M05, 68Q85}
\keywords{Moore flow, flow, reparametrization category with cuts,
	directed homotopy theory, enriched semicategory, semimonoidal category,
	combinatorial model category, left proper model category, Reedy category}


\begin{document}
	
	
	
\begin{abstract}
We introduce the notion of a reparametrization category with cuts. For every such reparametrization category \(\Pcat\), we prove the tensor lemma, namely that the tensor product of two objectwise weak homotopy equivalences of \(\Pcat\)-spaces is a weak equivalence, and then the left properness of the q-model structure of \(\Pcat\)-flows.  Finally, we prove that the interval reparametrization categories \(\G\), \(\M\), as well as the final category \(\One\), have cuts.  The last example recovers the left properness of the q-model structure of ordinary flows.
\end{abstract}

\maketitle
\setcounter{tocdepth}{1}
\tableofcontents
\hypersetup{linkcolor = dark-red}


\section{Introduction}

\subsection*{Presentation}
Directed algebraic topology studies spaces in which motion is constrained by a preferred direction. In classical topology, a path may always be traversed backwards, and the fundamental groupoid records this reversibility. This is not appropriate for phenomena in which time, causality, or an order of execution cannot be reversed. In the most familiar path-based models, a directed space is a topological space together with specified directed paths, stable under concatenation and under nondecreasing changes of parameter; maps and homotopies are required to preserve that direction. A directed path from \(x\) to \(y\) need not therefore have a reverse from \(y\) to \(x\), and the fundamental invariant is naturally a category rather than a groupoid. Grandis develops this point of view, as well as several other models of non-reversible phenomena, in \cite[Introduction and Part~I]{Grandis}.

Concurrency is a guiding source of examples. A point of a state space represents a global configuration of a concurrent system, and a directed path represents a possible execution. Independent actions are represented by higher-dimensional cubes: different directed routes across the same cube encode different interleavings of actions that may occur concurrently. Forbidden regions can encode, for example, violations of mutual exclusion. One consequently studies directed paths and their directed deformations, rather than merely the underlying undirected space. Spaces of such paths can distinguish essential schedules and detect phenomena such as deadlocks, unreachable states, and serializability. The survey \cite{GoubaultUserGuide} and the article \cite{FajstrupRaussenGoubault} introduce these ideas through concrete concurrent systems; the monograph \cite{DAT_book} gives a broader account of directed topology and concurrency. Together with \cite{Grandis}, these references provide entry points to the subject that do not presuppose the formalism used in this paper.

The formalism of flows records the part of this geometry that is relevant to execution. An ordinary flow \(X\) consists of a set \(X^0\) of states and, for each ordered pair \((\alpha,\beta)\), a topological space \(\mathbb P_{\alpha,\beta}X\) of nonconstant execution paths from \(\alpha\) to \(\beta\), together with associative composition of composable paths. Identity paths are omitted, so a flow is a topologically enriched semicategory. This retains more than the reachability relation: the homotopy type of \(\mathbb P_{\alpha,\beta}X\) records continuous families of executions from \(\alpha\) to \(\beta\).

There is a point-set issue hidden in path composition. If every path is parametrized by \([0,1]\), concatenating two paths requires rescaling them onto the two halves of that interval, and repeated binary concatenation is not strictly associative. The Moore construction remedies this by retaining a positive length. A path of length \(a\) followed by one of length \(b\) is concatenated on an interval of length \(a+b\); addition of lengths and blockwise concatenation are strictly associative. A reparametrization category \(\Pcat\) organizes both the allowed lengths and the allowed changes of parameter. A \(\Pcat\)-space is an enriched presheaf on \(\Pcat\), and its tensor product \(D\otimes E\) represents a first segment in \(D\) followed by a second segment in \(E\), modulo the allowed reparametrizations. A \(\Pcat\)-flow, or Moore flow, is a semicategory enriched in these \(\Pcat\)-spaces. The categories \(\G\) and \(\M\) have the positive real numbers as objects and addition as tensor product on objects: \(\G\) uses increasing homeomorphisms of intervals, whereas \(\M\) uses nondecreasing continuous surjections and therefore also allows pauses. This construction and its q-model structure are developed in \cite[Sections~4--8]{Moore1}. Thus Moore flows make concatenation strictly associative while retaining the dependence of execution paths on duration and reparametrization.

This article is a companion to \cite{leftproperflow}. That paper proves left properness for ordinary flows, where the operation combining spaces of successive paths is the cartesian product. The present paper identifies the additional structure that permits the same strategy for Moore flows and replaces the binary product throughout the cellular argument by the tensor product of \(\Pcat\)-spaces. This replacement is not formal: the tensor is defined by a coend, and its homotopical behavior must first be established for arbitrary, not necessarily cofibrant, \(\Pcat\)-spaces. When \(\Pcat\) is the final category, its tensor product is the cartesian product, a \(\Pcat\)-flow is an ordinary flow, and the main theorem of the present paper specializes to \cite[Theorem~5.6]{leftproperflow}.

Recall that a model category is \emph{left proper} if the pushout of a weak equivalence along a cofibration is again a weak equivalence \cite[Section~1.1]{MR99h:55031}. This is the basic homotopy-invariance statement for gluing along cofibrations: replacing an object by a weakly equivalent one before attaching a cell does not change the weak homotopy type after the attachment. It is therefore important when ordinary pushouts and colimits are used to model their derived, or homotopy, counterparts; for background on homotopy colimits, see \cite[Chapters~5--6]{Riehl}. Left properness is also a standard hypothesis in existence theorems for left Bousfield localizations \cite[Chapter~4]{ref_model2}.

For Moore flows, however, weak equivalences are detected objectwise on the path \(\Pcat\)-spaces, whereas colimits of Moore flows are not computed objectwise on those path objects. The reason is fundamental to directed composition: \emph{colimits can create execution paths}. Attaching one globe not only adds the paths in the attached cell; closure under composition also adds all finite composites in which a new path is interspersed with old ones. Likewise, identifying two states can make formerly noncomposable paths composable and thereby create reduced words of arbitrary finite length. Consequently, left properness of spaces or even of \(\Pcat\)-spaces does not by itself imply left properness of \(\Pcat\)-flows. Section~\ref{sec:6} controls these new words by cubical and Reedy latching objects, iterated tensor pushout products, and transfinite homotopy-colimit arguments. The tensor lemma proved in Section~\ref{sec:5} is indispensable here: the factors representing old execution paths need not be cofibrant, so the proof requires the tensor product to preserve weak equivalences without any cofibrancy assumption.

\subsection*{Main results.}
Throughout, \(\Top\) is either the category of \(\Delta\)-generated spaces or the category of \(\Delta\)-Hausdorff \(\Delta\)-generated spaces. The main results are the following.
\begin{enumerate}
	\item
	A \emph{cut structure} on a reparametrization category \(\Pcat\) consists, in particular, of continuous block-equivariant operators that cut a map \(L\to a\otimes b\) into maps \(L\to a\) and \(L\to b\), together with a block-equivariant homotopy from the original map to the map reconstructed from its two cuts. For every reparametrization category with cuts, every pair of arbitrary \(\Pcat\)-spaces \(D,E\), and every object \(L\), Theorem~\ref{thm:common-comparison} constructs natural maps
	\[
	r_L:(D\otimes E)(L)\longrightarrow D(L)\times E(L),
	\qquad
	s_L: D(L)\times E(L)\longrightarrow(D\otimes E)(L),
	\]
	and proves that they are homotopy inverses. The maps and the homotopies are natural in \(D\) and \(E\), and no cofibrancy hypothesis is imposed.
	
	\item
	Theorem~\ref{thm:tensor}, the tensor lemma, states that if \(f: D\to D'\) and \(g: E\to E'\) are objectwise weak homotopy equivalences of arbitrary \(\Pcat\)-spaces, then
	\[
	f\otimes g: D\otimes E\longrightarrow D'\otimes E'
	\]
	is an objectwise weak homotopy equivalence.
	
	\item
	Theorem~\ref{thm:left-proper} proves that, for every reparametrization category \(\Pcat\) with cuts, the q-model structure on \(\Pcat\)-flows is left proper.
	
	\item
	Theorem~\ref{thm:GM-cuttable} proves that both interval reparametrization categories \(\G\) and \(\M\) have cuts. Corollary~\ref{cor:GM-left-proper} therefore gives the left properness of the q-model structures of \(\G\)-flows and \(\M\)-flows. Proposition~\ref{prop:one-cuttable} proves that the final category \(\mathbf 1\) has cuts, and Corollary~\ref{cor:one-left-proper} recovers the left properness of ordinary flows.
	
	\item
	Proposition~\ref{prop:GM-nonunital} proves that neither \(\G\) nor \(\M\) is unital; more strongly, there is no object \(e\) and no natural isomorphism \(e\otimes L\cong L\). Theorem~\ref{thm:example} then shows that cuts are compatible with symmetry and strict unitality in a genuinely nontrivial example: the one-object enriched category \(\Pcat_{\times}\) whose endomorphism space is \([0,1]\) and whose composition and tensor product are multiplication is a symmetric strictly unital reparametrization category with cuts, but is not even equivalent to the final category as an ordinary category. Theorem~\ref{thm:left-proper} therefore also applies to the q-model structure of \(\Pcat_{\times}\)-flows.
\end{enumerate}

\subsection*{Organization of the paper.}
Section~\ref{sec:2} fixes the two convenient categories of spaces and recalls the q- and h-model structures and relative-\(T_1\) inclusions. Section~\ref{sec:3} recalls reparametrization categories, \(\Pcat\)-spaces and their tensor product, Moore flows, and the q-model structure. Section~\ref{sec:4} introduces reparametrization categories with cuts. Section~\ref{sec:5} proves the comparison theorem and deduces the tensor lemma. Section~\ref{sec:6} develops the cellular machinery needed for colimits of Moore flows---tensor pushout products, the free-cell latching analysis, globe attachments, and state identifications---and proves left properness by transfinite induction and a retract argument. Section~\ref{sec:7} constructs cuts for \(\G\) and \(\M\). Section~\ref{sec:8} treats the final category and explains how the theorem for ordinary flows is recovered. Finally, Section~\ref{sec:9} proves that \(\G\) and \(\M\) are not unital and constructs the symmetric strictly unital non-final example \(\Pcat_{\times}\).

\subsection*{Warning}

This paper is not completely self-contained, since it builds on \cite{leftproperflow}. Much of the argument adapts that work by replacing topological spaces with \(\Pcat\)-spaces and the cartesian product with the tensor product of \(\Pcat\)-spaces. Nevertheless, apart from a few specific technical points, the present paper can be read independently of \cite{leftproperflow}.

\section{Topological and model-categorical conventions}
\label{sec:2}

Throughout the paper, \(\Top\) denotes either the category of \(\Delta\)-generated spaces or its full subcategory of \(\Delta\)-Hausdorff \(\Delta\)-generated spaces.  A \(\Delta\)-generated space is a space whose topology is final with respect to all continuous maps from the standard topological simplices.  A \(\Delta\)-generated space is \(\Delta\)-Hausdorff if the image of every continuous map \([0,1]\to X\) is closed.  These definitions and both choices of \(\Top\) are studied in \cite[Section~2]{Moore1} and \cite[Section~2 and Appendix~B]{leftproperflow}.  In either case \(\Top\) is complete, cocomplete, locally presentable, and cartesian closed.  Its internal mapping space \(\TOP(-,-)\) is obtained by \(\Delta\)-kelleyfying the compact-open topology \cite[Section~2, pp.~3--4]{Moore1}.

We use two model structures on \(\Top\).  In the q-model structure, the weak equivalences are the weak homotopy equivalences and the fibrations are the Serre fibrations.  Here a weak homotopy equivalence means a map inducing a bijection on path components and isomorphisms on all based homotopy groups in positive degrees. For either choice of \(\Top\), its canonical self-enrichment is topologically bicomplete, with tensor \(X\otimes K=X\times K\) and cotensor \(X^{K}=\TOP(K,X)\); compare \cite[Definitions~1.4--1.5]{SchwanzlVogt}.  Local presentability implies the monomorphism hypothesis by \cite[Remark~5.20]{Barthel-Riel}, and hence the h-model structure exists by \cite[Corollary~5.23]{Barthel-Riel}. Its weak equivalences are the homotopy equivalences, and its cofibrations, called h-cofibrations, are the strong cofibrations of \cite[Definition~5.3]{Barthel-Riel}. Every q-cofibration is an h-cofibration; see, e.g., \cite[Remark~2.1]{leftproperflow}.  

A one-to-one map \(i: A\to X\) is a \emph{relative-\(T_1\) inclusion} if, for every open \(U\subset A\) and every \(z\in X\setminus i(U)\), there is an open \(W\subset X\) such that \(i(U)\subset W\) and \(z\notin W\) \cite[p. 686]{hocolimfacile}.  Every h-cofibration is relative-\(T_1\) by \cite[Definition~2.2 and Proposition~2.6]{leftproperflow}.

A \textit{transfinite tower} \(X:\lambda\to\mathcal C\) is a functor
whose canonical lift
\(
\widetilde X:\lambda\longrightarrow (X(0)\downarrow\mathcal C)
\)
is colimit-preserving.


\section{Reparametrization categories and Moore flows}
\label{sec:3}


\begin{definition}[reparametrization category]\label{def:reparametrization}
A \emph{reparametrization category} \((\Pcat,\otimes)\) is a small \(\Top\)-enriched semimonoidal category satisfying the following axioms.
\begin{enumerate}
\item The semimonoidal structure is strict; in particular, \(\otimes\) induces a semigroup structure on \(\Obj(\Pcat)\).
\item Every enriched hom-space \(\Pcat(\ell,\ell')\) is contractible.
\item If \(\phi:\ell\to\ell'\) and \(\ell'=\ell'_1\otimes\ell'_2\), there are objects \(\ell_1,\ell_2\) and maps \(\phi_i:\ell_i\to\ell'_i\) such that \(\ell=\ell_1\otimes\ell_2\) and \(\phi=\phi_1\otimes\phi_2\).
\end{enumerate}
The continuous map 
\[
\otimes:\Pcat(\ell_1,\ell_2) \times \Pcat(\ell'_1,\ell'_2) \longrightarrow \Pcat(\ell_1\otimes \ell'_1,\ell_2\otimes \ell'_2)
\]
is sometimes called the \textit{block tensor} in this paper.
\end{definition}


This is the axiomatization of \cite[Definition~4.3 and Appendix~A]{Moore1}; it is also recalled in \cite[Definition~1]{Moore3}.  A \(\Pcat\)-space is an enriched functor \(D:\Pcat^{\mathrm{op}}\to\Top\).  Thus a map
\(\alpha: a\to a'\) induces a contravariant action \(D(\alpha): D(a')\to D(a)\), and the adjoint action map \(\Pcat(a,a')\times D(a')\to D(a)\) is continuous.  The free \(\Pcat\)-space generated by \(U\in\Top\) at \(\ell\in\Obj(\Pcat)\) is
\begin{equation}\label{eq:free}
 \Ffree_\ell U=\Pcat(-,\ell)\times U;
\end{equation}
see \cite[Notation~5.15]{Moore1}.

The category \(\PSp\) has the projective q-model structure.  Its weak equivalences and fibrations are objectwise, and a set of generating q-cofibrations is
\begin{equation}\label{eq:Pspace-generators}
 \Ffree_\ell \mathsf{S}^{n-1}\longrightarrow \Ffree_\ell \mathsf{D}^n
 \qquad(\ell\in\Obj(\Pcat),\ n\geq0),
\end{equation}
where \(\mathsf{S}^{-1}=\varnothing\) and \(\mathsf{D}^0=\{0\}\); see \cite[Theorem~6.2]{dgrtop} and \cite[Section~8]{Moore1}.  This projective model structure is proper: Definition~\ref{def:reparametrization}(2) makes \(\Pcat\) locally contractible in the terminology of \cite{dgrtop}, and properness follows from \cite[Corollary~6.5]{dgrtop}.

For two \(\Pcat\)-spaces, their semimonoidal product is
\begin{equation}\label{eq:coend}
 (D\otimes E)(L) =\int^{a,b\in\Obj(\Pcat)}
   \Pcat(L,a\otimes b)\times D(a)\times E(b).
\end{equation}
This formula and the fact that it yields a biclosed semimonoidal structure are \cite[Definition~5.7 and Theorem~5.14]{Moore1}.  Write \([\psi,x,y]\) for the coend class of \((\psi,x,y)\).  For maps \(\alpha: a\to a'\) and \(\beta: b\to b'\), the defining relation is
\begin{equation}\label{eq:relation}
 [\,(\alpha\otimes\beta)\psi,x',y'\,]
 =[\,\psi,D(\alpha)x',E(\beta)y'\,].
\end{equation}
This quotient description is \cite[Corollary~5.13]{Moore1} for \(n=2\).

\begin{definition}\label{def:Pflow}
A \textit{\(\Pcat\)-flow} is a small semicategory enriched in the biclosed semimonoidal category \((\PSp,\otimes)\).  Thus it consists of a set \(X^0\) of states, a \(\Pcat\)-space \(\Path_{\alpha,\beta}X\) for every ordered pair of states, and associative enriched composition maps
\[
 \Path_{\alpha,\beta}X\otimes\Path_{\beta,\gamma}X
 \longrightarrow\Path_{\alpha,\gamma}X.
\]
Its \textit{total path \(\Pcat\)-space} is \(\Path X=\coprod_{(\alpha,\beta)\in X^0\times X^0} \Path_{\alpha,\beta}X\). For a \(\Pcat\)-space \(D\), the \textit{globe} \(\operatorname{Glob}(D)\) has states \(0,1\), path object \(D\) from \(0\) to \(1\), all other path objects empty, and no nontrivial composition.
\end{definition}

These definitions are \cite[Definitions~6.1, 6.2 and~6.5]{Moore1}.  The q-model structure on \(\PFlow\) is cofibrantly generated.  A map \(f: X\to Y\) of \(\PFlow\) is a weak equivalence precisely when \(f^0\) is a bijection and every map
\[
 \Path_{\alpha,\beta}X\longrightarrow
 \Path_{f(\alpha),f(\beta)}Y
\]
is an objectwise weak homotopy equivalence.  Its generating q-cofibrations are
\begin{equation}\label{eq:Pflow-generators}
\begin{split}
 C&:\varnothing\longrightarrow\{0\},\\
 R&:\{0,1\}\longrightarrow\{0\},\\
 \operatorname{Glob}(\Ffree_\ell \mathsf{S}^{n-1})
 &\longrightarrow \operatorname{Glob}(\Ffree_\ell \mathsf{D}^n)
 \qquad(\ell\in\Obj(\Pcat),\ n\geq0).
\end{split}
\end{equation}
The \(\Pcat\)-flows in the first two lines have no paths.  These assertions are \cite[Notation~8.2, Theorems~8.8 and~8.16]{Moore1}.  The convention \(n=0\) includes \(\varnothing=\mathsf{S}^{-1}\to \mathsf{D}^0\).

\section{Reparametrization categories with cuts}
\label{sec:4}

\begin{definition}[reparametrization category with cuts]
\label{def:cuttable}
A \emph{cut structure} on a reparametrization category \(\Pcat\) consists of the following data.
\begin{itemize}
\item For every object \(L\), a map
\[
 \delta_L: L\longrightarrow L\otimes L.
\]
\item For all \(L,a,b\in\Obj(\Pcat)\), continuous maps
\[
 c^{a,b}_1:\Pcat(L,a\otimes b)\longrightarrow\Pcat(L,a),
 \qquad
 c^{a,b}_2:\Pcat(L,a\otimes b)\longrightarrow\Pcat(L,b),
\]
called the \emph{cut operators}.  They are \textit{block-equivariant} in the following sense: for
\(\alpha: a\to a'\), \(\beta: b\to b'\), and
\(\psi: L\to a\otimes b\),
\begin{align}
 c^{a',b'}_1((\alpha\otimes\beta)\psi)
   &=\alpha c^{a,b}_1(\psi),                         \label{eq:cut-equiv1}\\
 c^{a',b'}_2((\alpha\otimes\beta)\psi)
   &=\beta c^{a,b}_2(\psi).                         \label{eq:cut-equiv2}
\end{align}
\item Put
\begin{equation}\label{eq:normalization}
 N^{a,b}_L(\psi)=
 \bigl(c^{a,b}_1(\psi)\otimes c^{a,b}_2(\psi)\bigr)\delta_L.
\end{equation}
There is a continuous homotopy
\[
 \Theta^{a,b}_L:
 [0,1]\times\Pcat(L,a\otimes b)\longrightarrow\Pcat(L,a\otimes b)
\]
from the identity to \(N^{a,b}_L\), and it is block-equivariant:
\begin{align}
 &\Theta^{a,b}_L(0,\psi)=\psi,\qquad
 \Theta^{a,b}_L(1,\psi)=N^{a,b}_L(\psi),
                                                        \label{eq:Theta-ends}\\
 &\Theta^{a',b'}_L(u,(\alpha\otimes\beta)\psi)
   =(\alpha\otimes\beta)\Theta^{a,b}_L(u,\psi).
                                                        \label{eq:Theta-equiv}
\end{align}
\end{itemize}
A \textit{reparametrization category with cuts} is a reparametrization category which admits a cut structure. We also say that a reparametrization category \textit{has cuts}.
\end{definition}

We prove that the interval reparametrization categories \(\G\) and \(\M\) have cuts in Theorem~\ref{thm:GM-cuttable}. We prove that the final category is a reparametrization category with cuts in Proposition~\ref{prop:one-cuttable}. Finally, we exhibit an example of a symmetric strictly unital reparametrization category with cuts which is not equivalent to the final category in Theorem~\ref{thm:example}.
 

\section{The tensor lemma}
\label{sec:5}


Fix a reparametrization category \(\Pcat\) with cuts, two arbitrary \(\Pcat\)-spaces \(D,E\), and \(L\in\Obj(\Pcat)\).  Define
\begin{equation}\label{eq:common-r}
 r_L:(D\otimes E)(L)\longrightarrow D(L)\times E(L)
\end{equation}
by
\begin{equation}\label{eq:common-r-formula}
 r_L[\psi,x,y]=
 \bigl(D(c^{a,b}_1(\psi))x,E(c^{a,b}_2(\psi))y\bigr).
\end{equation}
Equations \eqref{eq:cut-equiv1}--\eqref{eq:cut-equiv2}, contravariant functoriality, and \eqref{eq:relation} show that this is independent of the representative.  Continuity follows from continuity of the cut operators and of the enriched actions, together with the quotient universal property of the coend. The formula also proves naturality in \(D\) and \(E\).

Using the \((a,b)=(L,L)\) coend summand, define
\begin{equation}\label{eq:common-s}
 s_L: D(L)\times E(L)\longrightarrow(D\otimes E)(L),\qquad
 s_L(x,y)=[\delta_L,x,y].
\end{equation}
It is continuous because the coprojection of a coend summand is continuous.

\begin{thm}[comparison theorem]\label{thm:common-comparison}
For every reparametrization category with cuts \(\Pcat\), arbitrary \(\Pcat\)-spaces \(D,E\), and every \(L\in\Obj(\Pcat)\), the maps \(r_L\) and \(s_L\) are homotopy inverses.  Thus, naturally in \(D\) and \(E\),
\[
 (D\otimes E)(L)\simeq D(L)\times E(L).
\]
No cofibrancy hypothesis is required.
\end{thm}

\begin{proof}
Put \(\lambda_i=c^{L,L}_i(\delta_L)\in\Pcat(L,L)\).  Then
\[
 r_Ls_L(x,y)=\bigl(D(\lambda_1)x,E(\lambda_2)y\bigr).
\]
The space \(\Pcat(L,L)\) is contractible by
Definition~\ref{def:reparametrization}(2), hence it is path connected.  To
spell this out, a contraction to a point gives paths from \(\lambda_i\) and
\(\id_L\) to that point; concatenating the first with the reverse of the
second gives a path \(\lambda_{i,u}\) from \(\lambda_i\) to \(\id_L\).  Applying
the continuous enriched actions of \(D\) and \(E\) to these paths yields
\(r_Ls_L\simeq\id\).

For \(\psi\in\Pcat(L,a\otimes b)\), relation \eqref{eq:relation}, applied in
the \((L,L)\) summand, gives
\begin{equation}\label{eq:common-sr}
 s_Lr_L[\psi,x,y]=[N^{a,b}_L(\psi),x,y].
\end{equation}
By \eqref{eq:Theta-equiv}, the maps
\[
 (u,[\psi,x,y])\longmapsto
 [\Theta^{a,b}_L(u,\psi),x,y]
\]
respect all coend relations.  They therefore descend to a homotopy on \((D\otimes E)(L)\).  The descent is continuous because \([0,1]\times-\) preserves coproducts and coequalizers: in the cartesian closed category \(\Top\), it is left adjoint to the internal hom functor \(\operatorname{TOP}([0,1],-)\).  By \eqref{eq:Theta-ends} and \eqref{eq:common-sr}, its endpoints are the identity and \(s_Lr_L\). Consequently \(s_Lr_L\simeq\id\).  

The homotopy from \(r_Ls_L\) to the identity is natural in \(D\) and \(E\) by naturality of the enriched actions.  The homotopy from \(s_Lr_L\) to the identity leaves \(x\) and \(y\) unchanged, and is therefore natural in \(D\) and \(E\) as well.
\end{proof}

\begin{thm}[tensor lemma]\label{thm:tensor}
Let \(\Pcat\) be a reparametrization category with cuts.  If \(f: D\to D'\) and \(g: E\to E'\) are objectwise weak homotopy equivalences of arbitrary \(\Pcat\)-spaces, then
\[
 f\otimes g: D\otimes E\longrightarrow D'\otimes E'
\]
is an objectwise weak homotopy equivalence.
\end{thm}

\begin{proof}
For every \(L\in\Obj(\Pcat)\), naturality of \(r_L\) gives a commutative square
\[
\begin{tikzcd}
 (D\otimes E)(L) \arrow[r,"(f\otimes g)_L"] \arrow[d,"r_L"']
   & (D'\otimes E')(L) \arrow[d,"r'_L"] \\
 D(L)\times E(L) \arrow[r,"f_L\times g_L"']
   & D'(L)\times E'(L)
\end{tikzcd}
\]
The vertical maps are homotopy equivalences by Theorem~\ref{thm:common-comparison}, hence weak homotopy equivalences.  The bottom map is a weak homotopy equivalence: path components of a finite product are products of path components, and maps and homotopies into a product give natural isomorphisms
\[
 \pi_n(X\times Y,(x,y))\cong
 \pi_n(X,x)\times\pi_n(Y,y)\qquad(n\geq1).
\]
The two-out-of-three property, applied to the square, proves that \((f\otimes g)_L\) is a weak homotopy equivalence.  Since \(L\) was arbitrary, the result is objectwise.
\end{proof}


\section{Left properness}
\label{sec:6}

For maps \(f: A\to B\) and \(g: C\to D\) of \(\Pcat\)-spaces, their tensor pushout product is
\[
 f\mathbin{\square_\otimes}g:
 (A\otimes D)\mathop{\coprod}_{A\otimes C}(B\otimes C)
 \longrightarrow B\otimes D.
\]
The notation \(\square_\times\) will mean the analogous pushout product in \((\Top,\times)\). Iterated pushout products are associated from the left; by associativity, any other parenthesization is canonically isomorphic to this one. This is \cite[Definition~9.1]{Moore1}.

\begin{proposition}[cubical formula for the iterated tensor pushout product]\label{prop:calculpushout} 
Let \[0\leq i\leq p,\quad f_i:A_i\longrightarrow B_i\] be \(p+1\) maps of \(\mathcal{P}\)-spaces. Let \(S\subset
\{0,\dots,p\}\). Let 
\[C_p(S):=C_0 \otimes \dots \otimes C_p \hbox{ with } \begin{cases}
	C_i = A_i & \hbox{ if }i\notin S\\
	C_i = B_i & \hbox{ if }i\in S.
\end{cases}\]
If \(S\) and \(T\) are two subsets of \(\{0,\dots,p\}\) such that \(S\subset
T\), let \[C_p(i_S^T):C_p(S)\longrightarrow C_p(T)\] be the morphism 
\[g_0 \otimes \dots \otimes g_p \hbox{ with } \begin{cases}
	g_i = \id_{B_i} & \hbox{ if }i\in S\\
	g_i = f_i & \hbox{ if }i\in T\backslash S\\
	g_i = \id_{A_i} & \hbox{ if }i\notin T.
\end{cases}\]
Then:  
\begin{enumerate} 
	\item the assignments \(S\mapsto C_p(S)\) and \(i_S^T\mapsto C_p(i_S^T)\) define a functor
	\[
	\Pcat(\{0,\dots,p\})\longrightarrow
	\PSp.
	\]
	\item there exists a canonical morphism 
	\[\liminj_{S\subsetneqq \{0,\dots,p\}} 
	C_p(S)\longrightarrow C_p(\{0,\dots,p\}).\]
	and it is equal to the morphism \(f_0\square_\otimes\dots \square_\otimes f_p\). 
\end{enumerate}
\end{proposition}

\begin{proof}
	This is \cite[Proposition~9.8]{Moore1}.
\end{proof}


Proposition~\ref{prop:calculpushout} is called the \emph{cubical formula} for the iterated tensor pushout product because the simplices of the order complex of the chain \(\{0<\cdots<p\}\), including the empty simplex, form a poset naturally isomorphic to the Boolean lattice \(\{0<1\}^{p+1}\): a simplex is sent to its characteristic function. The objects and maps appearing in Proposition~\ref{prop:calculpushout} therefore assemble into a diagram of \(\Pcat\)-spaces
\[
D: \{0<1\}^{p+1}\longrightarrow \PSp,
\qquad
D(\epsilon_0,\ldots,\epsilon_p)
=
C_p\bigl(\{i\mid \epsilon_i=1\}\bigr).
\]
The indexing poset is a direct Reedy category with degree function
\[
d(\epsilon_0,\ldots,\epsilon_p)=\sum_{i=0}^{p}\epsilon_i.
\]
With respect to this Reedy structure, the iterated tensor pushout product
\[
f_0\square_\otimes\cdots\square_\otimes f_p
\]
is precisely the latching map of \(D\) at the terminal vertex
\((1,\ldots,1)\):
\[
L_{(1,\ldots,1)}D
=
\liminj_{
	\partial\bigl(
	\{0<1\}^{p+1}\mathbin{\downarrow}(1,\ldots,1)
	\bigr)
}D
\longrightarrow
D(1,\ldots,1).
\]

\begin{lem}\label{lem:finite-tensor-coproduct}
Nonempty finite tensor products and arbitrary coproducts of objectwise weak homotopy equivalences of \(\Pcat\)-spaces are objectwise weak homotopy equivalences.
\end{lem}

\begin{proof}
The assertion about finite tensor products follows by induction from Theorem~\ref{thm:tensor}; the case of one factor is tautological.  For coproducts, fix \(L\in\Obj(\Pcat)\).  A coproduct of spaces is a disjoint union. Its set of path components is the disjoint union of the sets of path components of its summands.  Moreover, a based map from the connected sphere \(\mathsf{S}^k\) (\(k\geq1\)), and every based homotopy between such maps, lies in the summand containing the basepoint.  Hence the homotopy groups of a coproduct, at a chosen basepoint, are the homotopy groups of that summand. The claimed assertion follows directly from the definition of weak homotopy equivalence.
\end{proof}

\begin{lem}[well known] \label{lem:h-cartesian-monoidal}
	The h-model structure on \(\Top\) is cartesian monoidal.
\end{lem}

\begin{proof}
	The usual interval is a good cylinder object since the inclusion map \(\{0,1\}\subset [0,1]\) is an h-cofibration, being a q-cofibration; see \cite[Example~1.11(A)(1) and Sections~4.1 and~4.3]{SchwanzlVogt}. By \cite[Definition~5.3 and Corollary~5.23]{Barthel-Riel}, the cofibrations of the h-model structure are the maps having the left lifting property with respect to the h-fibrations that are
	homotopy equivalences.  By \cite[Proposition~3.5(6)]{SchwanzlVogt}, these are precisely the strong cofibrations of \cite{SchwanzlVogt}.
	
	Let \(i: A\to B\) and \(j: C\to D\) be \(h\)-cofibrations. Applied to the closed two-variable adjunction determined by the cartesian product, \cite[Pairing Theorem~2.7(1)]{SchwanzlVogt} shows that
	\[
	i\square_{\times}j:
	(A\times D)\amalg_{A\times C}(B\times C)\longrightarrow B\times D
	\]
	is a strong cofibration.  If \(i\) or \(j\) is trivial, then it is a homotopy equivalence, and \cite[Addendum~3.6(1)]{SchwanzlVogt} shows that \(i\square_{\times}j\) is a homotopy equivalence.  Thus the cartesian product satisfies the pushout-product axiom; equivalently, it is a Quillen bifunctor by \cite[Lemma~4.2.2]{MR99h:55031}.
	
	Finally, every object is strongly cofibrant by \cite[Lemma~2.6(3)]{SchwanzlVogt}.  In particular, the cartesian unit \(\{*\}\) is cofibrant, so the unit axiom is automatic: one may take \(\id_{\{*\}}\) as its cofibrant replacement; see \cite[Definition~4.2.6]{MR99h:55031}.  Therefore the h-model structure on \(\Top\) is cartesian monoidal.
\end{proof}

\begin{lem}[well known]\label{lem:product-hcof}
If \(i: A\to B\) is an h-cofibration and \(K\) is any space, then \(K\times i: K\times A\to K\times B\) is an h-cofibration.  Pushouts and transfinite composites of h-cofibrations are h-cofibrations.
\end{lem}

\begin{proof}
The map \(K\times i\) is equal to the pushout product 
\(
(\varnothing\to K) \square_\times i
\). 
Since every space is h-cofibrant by \cite[Lemma~2.6(3)]{SchwanzlVogt}, Lemma~\ref{lem:h-cartesian-monoidal} implies that this map is an h-cofibration. The last assertion is a general feature of the class of cofibrations of a model category, here the h-model category of spaces. 
\end{proof}


The next lemma is the point-set substitute for the assertion, used for ordinary flows, that a product pushout product with arbitrary ordinary path spaces is an h-cofibration.

\begin{lem}[free-cell latching lemma]\label{lem:free-latching}
	Let \(p,k\geq0\) with \(p+k>0\), and let \(W_1,\dots,W_p\) be arbitrary \(\Pcat\)-spaces.  For \(1\leq r\leq k\), choose \(\ell_r\in\Obj(\Pcat)\) and \(n_r\geq0\), and let
	\[
	j_r=\Ffree_{\ell_r}
	(\mathsf{S}^{n_r-1}\hookrightarrow \mathsf{D}^{n_r})
	\]
	be a generating projective q-cofibration of \(\Pcat\)-spaces.  Every evaluation at \(L\in\Obj(\Pcat)\) of an iterated tensor pushout product made, in any fixed order, from the maps \(\varnothing\to W_i\) and the maps \(j_r\) is an h-cofibration.  The same conclusion holds if any of the maps \(j_r\) is replaced by a pushout of \(j_r\) in \(\PSp\).
\end{lem}

\begin{proof}
	We first record explicitly how the topological variables in free \(\Pcat\)-spaces pass through the tensor product.  For a space \(U\), the pointwise copower \(D\times U\) satisfies natural homeomorphisms
	\begin{align}
		((D\times U)\otimes E)(L)&\cong(D\otimes E)(L)\times U,\label{eq:copower-left}\\
		(D\otimes(E\times U))(L)&\cong(D\otimes E)(L)\times U.\label{eq:copower-right}
	\end{align}
	Indeed, in the coend formula \eqref{eq:coend}, product with \(U\) commutes with the coproduct and coequalizer defining the coend: the functor \(-\times U\) is a left adjoint because \(\Top\) is cartesian closed.  The same argument in the other variable proves the second homeomorphism.  These homeomorphisms are natural in all their variables because they are induced by the canonical colimit comparison maps.  By \eqref{eq:free}, they apply to a free \(\Pcat\)-space through the identity
	\begin{equation}\label{eq:free-copower}
		\Ffree_\ell U=(\Ffree_\ell\{*\})\times U.
	\end{equation}
	
	We now compute the latching map rather than only its codomain.  Number the \(p+k\) maps in their prescribed order, and put \(I=\{1,\ldots,p+k\}\).  Let \(I_W\subset I\) be the positions occupied by the maps \(\varnothing\to W_i\), and let \(I_F=I\setminus I_W\) be the positions occupied by the free cells.  For \(S\subset I\), let \(C(S)\) be the ordered tensor word obtained by choosing the target of the map in position \(q\) when \(q\in S\) and its source when \(q\notin S\).  The cubical formula for an iterated pushout product gives the canonical map
	\begin{equation}\label{eq:punctured-cube}
		\liminj_{S\subsetneq I}C(S)\longrightarrow C(I);
	\end{equation}
	see Proposition~\ref{prop:calculpushout}.
	
	The tensor product of \(\Pcat\)-spaces preserves colimits separately since it is a biclosed semimonoidal structure; in particular, a tensor word with an empty factor is empty.  It follows that \(C(S)=\varnothing\) unless \(I_W\subset S\).  Empty vertices contribute neither points nor relations to the usual coproduct--coequalizer construction of the colimit in \eqref{eq:punctured-cube}.  Thus, if \(k>0\), its source is the colimit of the subcube
	\[
	C(I_W\cup T),\qquad T\subsetneq I_F.
	\]
	
	Let \(K\) be the ordered tensor word obtained from \(C(I)\) by replacing every free factor \(\Ffree_{\ell_r}\mathsf{D}^{n_r}\) by \(\Ffree_{\ell_r}\{*\}\) and leaving every \(W_i\) unchanged.  For \(T\subset I_F\), relabeling the free cells so that their indices record their order of occurrence, set
	\[
	U_r(T)=
	\begin{cases}
		\mathsf{D}^{n_r},&\text{if the position of \(j_r\) belongs to \(T\)},\\
		\mathsf{S}^{n_r-1},&\text{otherwise}.
	\end{cases}
	\]
	Repeated use of \eqref{eq:copower-left}, \eqref{eq:copower-right}, and \eqref{eq:free-copower} gives a natural homeomorphism
	\begin{equation}\label{eq:free-cube-vertices}
		C(I_W\cup T)(L)\cong
		K(L)\times\prod_{r=1}^{k}U_r(T).
	\end{equation}
	Naturality in the spaces \(U_r(T)\) shows that, under \eqref{eq:free-cube-vertices}, every structure map of the cube is the identity of \(K(L)\) times the corresponding product of identities and inclusions \(\mathsf{S}^{n_r-1}\to \mathsf{D}^{n_r}\).  Colimits of enriched presheaves are computed objectwise \cite[Proposition~5.3]{dgrtop}, and \(K(L)\times-\) preserves colimits because \(\Top\) is cartesian closed. Consequently, evaluation of the whole map \eqref{eq:punctured-cube}, not merely of its terminal vertex, is canonically isomorphic to
	\begin{equation}\label{eq:free-latching-product}
		K(L)\times
		\bigl(i_1\mathbin{\square_\times}\cdots
		\mathbin{\square_\times}i_k\bigr),
		\qquad i_r:\mathsf{S}^{n_r-1}\hookrightarrow \mathsf{D}^{n_r},
	\end{equation}
	where the free cells are listed in their order of occurrence in the tensor word. The iterated pushout product in \eqref{eq:free-latching-product} is a CW-subcomplex inclusion---the union of the coordinate boundaries---and hence a q-cofibration since the q-model structure of k-spaces is cartesian monoidal by \cite[Proposition~4.2.11]{MR99h:55031}. The source and the target of this q-cofibration are Hausdorff \(\Delta\)-generated spaces. Thus the map is a q-cofibration of \(\Top\). It is therefore an h-cofibration of \(\Top\), and its product with \(K(L)\) is an h-cofibration by Lemma~\ref{lem:product-hcof}.  If \(k=0\), every \(C(S)\) for \(S\subsetneq I\) in \eqref{eq:punctured-cube} is empty, so the evaluated map is \(\varnothing\to C(I)(L)\); this is an h-cofibration because every space is h-cofibrant.
	
	Finally, the functor that forms a pushout product preserves pushouts in each arrow variable: this follows by writing its source as a pushout and using the separate colimit preservation of \(\otimes\).  Replacing some \(j_r\) by pushouts of \(j_r\) therefore replaces the map just treated by a finite composite of pushouts of such maps.  It remains an h-cofibration by Lemma~\ref{lem:product-hcof}.
\end{proof}

\begin{lem}[Reedy comparison]\label{lem:Reedy-comparison}
Let \(\mathcal R\) be a small Reedy category for which the colimit functor is left Quillen.  Let \(D,D':\mathcal R\to\Top\) be Reedy cofibrant for the h-model structure.  If \(D\to D'\) is objectwise a weak homotopy equivalence, then
\[
 \liminj_{\mathcal R}D\longrightarrow\liminj_{\mathcal R}D'
\]
is a weak homotopy equivalence.
\end{lem}

\begin{proof}
Because \(D\) and \(D'\) are Reedy h-cofibrant and colimit is left Quillen, their ordinary colimits compute their h-model homotopy colimits.  For every small diagram of spaces, its q-model and h-model homotopy colimits have the same weak homotopy type by the Dugger--Isaksen comparison (\cite[Theorem~A.7]{hocolimfacile} adapted for \(\Delta\)-generated spaces in \cite[Theorem~2.8]{leftproperflow}). The q-model homotopy colimit, being a derived functor, carries an objectwise q-weak equivalence to a q-weak equivalence.  Applying these facts to the map \(D\to D'\) and using the two-out-of-three property gives the result.  This is precisely the homotopy-colimit step of \cite[proof of Theorem~5.3]{leftproperflow}.
\end{proof}

\begin{cor}[transfinite gluing]\label{cor:transfinite-gluing}
Let \(X\to Y\) be a map of transfinite towers \(\lambda\to\Top\).  Suppose that all successor maps in both towers are h-cofibrations and that \(X_\mu\to Y_\mu\) is a weak homotopy equivalence for every \(\mu<\lambda\). Then \(\liminj X\to\liminj Y\) is a weak homotopy equivalence.
\end{cor}

\begin{proof}
	Regard the ordinal \(\lambda\) as a direct category, and hence as a Reedy category, with degree function \(d(\mu)=\mu\). For either tower, the latching map at a successor ordinal \(\mu+1\) is \(X_\mu\to X_{\mu+1}\); at a nonzero limit ordinal \(\mu\), it is the canonical map
	\[
	\liminj_{\nu<\mu}X_\nu\longrightarrow X_\mu,
	\]
	which is an isomorphism by definition of a transfinite tower; and at \(0\), it is \(\varnothing\to X_0\). The successor latching maps are h-cofibrations by hypothesis, the limit latching maps are isomorphisms, and the latching map at \(0\) is an h-cofibration because every space is h-cofibrant. Thus both towers are Reedy h-cofibrant. Moreover, the constant-diagram functor, which is right adjoint to the colimit functor, preserves h-fibrations and trivial h-fibrations because these are defined objectwise. Consequently, the colimit functor is left Quillen; see \cite[Definition~5.1.1 and Corollary~5.1.6]{MR99h:55031}. Lemma~\ref{lem:Reedy-comparison} therefore applies.
\end{proof}


We recall the Reedy description needed below.  Given a set \(S\) and \(u,v\in S\), let \(\mathcal R_{u,v}(S)\) denote the category written \(\Pcat_{u,v}(S)\) in \cite[Section~9]{Moore1}; the change of letter avoids confusion with the reparametrization category \(\Pcat\).  Its objects are finite words
\[
 ((u_0,\epsilon_1,u_1),\ldots,
   (u_{m-1},\epsilon_m,u_m)),
\]
where \(u_i\in S\), \(\epsilon_i\in\{0,1\}\), and \(\epsilon_i=1\) forces \((u_{i-1},u_i)=(u,v)\).  Its arrows are generated by contractions of two adjacent \(0\)-labelled entries (the composition maps) and by replacements \((u,0,v)\to(u,1,v)\) (the inclusion maps), subject to the associativity and interchange relations.  The complete presentation is \cite[Section~9, before Notation~9.5]{Moore1}.  It is a Reedy category with degree \(m+\sum_i\epsilon_i\); inclusion maps raise degree and composition maps lower it.  For a diagram \(D\), the latching object at
\(\mathbf n\) is
\[
 L_{\mathbf n}D=
 \liminj_{\partial(\mathcal R_{u,v}(S)^+\downarrow\mathbf n)}D.
\]
The Reedy model structure exists for every model target, and its colimit functor is left Quillen \cite[Theorem~9.6]{Moore1}.

Consider a pushout of \(\Pcat\)-flows
\begin{equation}\label{eq:globe-pushout}
\begin{tikzcd}
 \operatorname{Glob}(\partial Z) \arrow[r,"g"] \arrow[d]
   & A \arrow[d,"f"] \\
 \operatorname{Glob}(Z) \arrow[r] & X \cocartesian
\end{tikzcd}
\end{equation}
and put \(u=g(0)\), \(v=g(1)\).  Define
\begin{equation}\label{eq:T-pushout}
 T=Z\amalg_{\partial Z}\Path_{u,v}A.
\end{equation}
The diagram
\(D_f:\mathcal R_{u,v}(A^0)\to\PSp\) assigns to the word
\(\mathbf n\) above the ordered tensor product
\begin{equation}\label{eq:Df}
 D_f(\mathbf n)=Z_1\otimes\cdots\otimes Z_m,
 \qquad
 Z_i=\begin{cases}
  \Path_{u_{i-1},u_i}A,&\epsilon_i=0,\\
  T,&\epsilon_i=1.
 \end{cases}
\end{equation}
The inclusion arrows use \(\Path_{u,v}A\to T\) and the contraction arrows use composition in \(A\).  This is well defined and
\begin{equation}\label{eq:colim-Df}
 \liminj D_f\cong\Path X
\end{equation}
by \cite[Theorem~9.7]{Moore1}.  Its latching map is the iterated tensor pushout product of the maps \(\varnothing\to\Path_{a,b}A\) belonging to the \(0\)-entries and the maps \(\Path_{u,v}A\to T\) belonging to the \(1\)-entries; if all entries are \(0\), the latching object is empty.  This is \cite[Proposition~9.9]{Moore1}. Colimits of enriched presheaves are computed objectwise \cite[Proposition~5.3]{dgrtop}; consequently, evaluation at \(L\) commutes with the latching colimits and with the colimit in \eqref{eq:colim-Df}.

\begin{proposition}[globe step]\label{prop:globe-step}
Let \(s: A\to A'\) be a weak equivalence of \(\Pcat\)-flows.  Let \(\partial Z\to Z\) be one of the generating maps \(\Ffree_\ell \mathsf{S}^{n-1}\to\Ffree_\ell \mathsf{D}^n\), and form the two pushouts
\[
 X=A\mathop{\coprod}_{\operatorname{Glob}(\partial Z)}
       \operatorname{Glob}(Z),
 \qquad
 X'=A'\mathop{\coprod}_{\operatorname{Glob}(\partial Z)}
       \operatorname{Glob}(Z),
\]
where the second attaching map is the first one followed by \(s\).  Then the induced map \(\bar s: X\to X'\) is a weak equivalence.  Moreover \(\Path A(L)\to\Path X(L)\) is an h-cofibration for every \(L\in\Obj(\Pcat)\).
\end{proposition}

\begin{proof}
Since \(s\) is a weak equivalence, it is bijective on states.  After relabelling \(A'^0\), we may suppose that both state sets and the attaching endpoints \(u,v\) agree.  Let
\[
 T=\Path_{u,v}A\mathop{\coprod}_{\partial Z}Z,
 \qquad
 T'=\Path_{u,v}A'\mathop{\coprod}_{\partial Z}Z.
\]
The map \(\partial Z\to Z\) is a projective q-cofibration of \(\Pcat\)-spaces, and the projective q-model structure on \(\PSp\) is left proper by \cite[Corollary~6.5]{dgrtop}.  Therefore \(T\to T'\) is an objectwise weak homotopy equivalence.

The map \(\Path_{u,v}A\to T\) is a pushout of the free cell \(\partial Z\to Z\).  Consequently, Proposition~9.9 of \cite{Moore1} and Lemma~\ref{lem:free-latching} imply that, after evaluation at any \(L\), all latching maps of \(D_f\) are h-cofibrations; the all-zero case uses \(\varnothing\to D_f(\mathbf n)(L)\).  Thus \(D_f(L)\) is Reedy h-cofibrant. The same reasoning applies to \(D_{f'}(L)\).

There is a natural map \(D_f\to D_{f'}\).  Every one of its vertices is a finite ordered tensor product of maps which are either \(\Path_{a,b}A\to\Path_{a,b}A'\) or \(T\to T'\).  All these factor maps are objectwise weak homotopy equivalences, so every vertex map is one by Lemma~\ref{lem:finite-tensor-coproduct}.  Lemma~\ref{lem:Reedy-comparison} and \eqref{eq:colim-Df} now give an objectwise weak homotopy equivalence
\[
 \Path X\cong\liminj D_f\longrightarrow
 \liminj D_{f'}\cong\Path X'.
\]
The globe attachment does not change the set of states: both globes have the same two states and the state map between them is the identity.  Hence \(X^0=A^0\to A'^0=X'^0\) is a bijection.  The map of total path spaces preserves the coproduct summand indexed by each ordered pair of states. Because these summands are open and closed, the calculation of \(\pi_0\) and the based homotopy groups used in Lemma~\ref{lem:finite-tensor-coproduct} shows that the restriction to every summand is a weak homotopy equivalence.  The characterization recalled after Definition~\ref{def:Pflow} therefore proves that \(\bar s\) is a weak equivalence.

It remains to prove the last assertion.  Use the map of diagrams \(D_{\id_A}\to D_f\) from the proof of \cite[Theorem~9.10]{Moore1}.  Its relative latching map is an isomorphism when all labels are zero; otherwise it is an iterated tensor pushout product of maps \(\varnothing\to\Path_{a,b}A\) and at least one copy of \(\Path_{u,v}A\to T\); see \cite[proof of Theorem~9.10, cases (a) and (b)]{Moore1}.  After evaluation at \(L\) it is an h-cofibration by Lemma~\ref{lem:free-latching}.  Thus \(D_{\id_A}(L)\to D_f(L)\) is a Reedy h-cofibration.  Since colimit is left Quillen \cite[Theorem~9.6]{Moore1}, its colimit \(\Path A(L)\to\Path X(L)\) is an h-cofibration.
\end{proof}


The following explicit calculation supplies both the weak-equivalence and the cofibration assertions for the generator \(R\).

\begin{lem}[reduced-word calculation]\label{lem:R-step}
Let \(a,b\in A^0\), and let \(X\) be the pushout of \(A\) along the map \(R:\{0,1\}\to\{0\}\) which sends \(0\) to \(a\) and \(1\) to \(b\). Then \(X^0=A^0/{\sim}\), where \(\sim\) identifies \(a\) and \(b\), and \(\Path A(L)\to\Path X(L)\) is an h-cofibration for every \(L\in\Obj(\Pcat)\).

If \(s: A\to A'\) is a weak equivalence and \(X'\) is the corresponding pushout from \(A'\), then the induced map \(X\to X'\) is a weak equivalence.
\end{lem}

\begin{proof}
If \(a=b\), the pushout is \(A\) and there is nothing to prove.  Suppose that \(a\neq b\).  Write \([c]\) for the class of a state \(c\in A^0\).  A \emph{reduced word} is a finite sequence
\[
 w=((\alpha_1,\beta_1),\ldots,(\alpha_m,\beta_m)),\qquad m\geq1,
\]
such that \([\beta_i]=[\alpha_{i+1}]\) but \(\beta_i\neq\alpha_{i+1}\) for \(1\leq i<m\).  Associate to it the \(\Pcat\)-space
\begin{equation}\label{eq:word-space}
 \Path_w A=
 \Path_{\alpha_1,\beta_1}A\otimes\cdots\otimes
 \Path_{\alpha_m,\beta_m}A.
\end{equation}
For two classes \(\bar\alpha,\bar\beta\in A^0/{\sim}\), put
\begin{equation}\label{eq:R-normal-form}
 \Path_{\bar\alpha,\bar\beta}X=
 \coprod_{\substack{w\ \mathrm{reduced}\\
 [\alpha_1]=\bar\alpha,\\
 [\beta_m]=\bar\beta}}
 \Path_w A.
\end{equation}

We justify that \eqref{eq:R-normal-form} is the pushout path object.  Two reduced words whose terminal and initial classes agree are composed by concatenation if their actual middle states are different, and, if those states agree, by applying the composition of \(A\) to the last factor of the first word and the first factor of the second.  The result is reduced: the two adjacent inequalities on either side, when present, are unchanged. These maps are maps of \(\Pcat\)-spaces because \(\otimes\) is functorial and the composition of \(A\) is enriched.  They are associative by the associativity of the composition of \(A\) (the only possible ambiguity is the order in which a consecutive block with equal actual intermediate states is composed).  Thus \eqref{eq:R-normal-form} defines a \(\Pcat\)-flow and the one-letter words define a map \(A\to X\) identifying \(a\) and \(b\).

If \(h: A\to Y\) is any map of \(\Pcat\)-flows with \(h(a)=h(b)\), there is a unique extension \(X\to Y\): on the summand \eqref{eq:word-space}, it is the finite composition in \(Y\) of the images of its factors.  This extension is continuous and enriched because it is a map on each coproduct summand, built from tensor products and the enriched composition maps of \(Y\). Uniqueness follows because every word is a composite of its one-letter factors.  This proves the pushout universal property and therefore \eqref{eq:R-normal-form}.

The one-letter summands in \eqref{eq:R-normal-form} form precisely \(\Path A\).  Hence \(\Path A(L)\to\Path X(L)\) is the inclusion of a collection of coproduct summands.  Such an inclusion is closed and has the homotopy extension property: on every complementary summand, extend a prescribed initial map by the constant homotopy.  It is therefore an h-cofibration.

Now let \(s: A\to A'\) be a weak equivalence.  Its bijection on states identifies the reduced-word indexing sets for \(A\) and \(A'\).  On the summand indexed by \(w\), the induced map is the finite tensor product of the objectwise weak homotopy equivalences on the factors of \eqref{eq:word-space}; it is an objectwise weak homotopy equivalence by Lemma~\ref{lem:finite-tensor-coproduct}.  Taking the coproduct over all words preserves this property by the same lemma.  The quotient map on states is a bijection because it is obtained from a bijection of state sets by identifying the corresponding pair.  Thus \(X\to X'\) is a weak equivalence.
\end{proof}


\begin{thm}\label{thm:left-proper}
For every reparametrization category with cuts \(\Pcat\), the q-model structure on \(\PFlow\) is left proper.
\end{thm}

\begin{proof}
First suppose that \(i: A\to X\) is a relative cell complex generated by the maps \eqref{eq:Pflow-generators}.  A pushout of a coproduct of generating maps may be refined, after well-ordering its cells, into a transfinite composite of pushouts of single generating maps: attach the cells one at a time using the composites of their original attaching maps, and use the universal property of coproducts and pushouts to identify the resulting colimit with the simultaneous pushout.  We may therefore write \(i\) as a transfinite tower indexed by \(\lambda+1\)
\[
 A=A_0\longrightarrow A_1\longrightarrow\cdots
 \longrightarrow A_\lambda=X
\]
in which each successor is a pushout of one map in \eqref{eq:Pflow-generators}.

Let \(s=s_0: A\to A'\) be a weak equivalence and base-change this tower along \(s\).  Thus
\[
 A'_\mu=A'\amalg_{A}A_\mu,
 \qquad s_\mu: A_\mu\longrightarrow A'_\mu,
\]
and \(A'_{\mu+1}\) is obtained from \(A'_\mu\) by the corresponding generating attachment.  We prove by transfinite induction that every \(s_\mu\) is a weak equivalence and, simultaneously, that every successor map
\begin{equation}\label{eq:successor-hcof}
 \Path A_\mu(L)\longrightarrow\Path A_{\mu+1}(L),
 \qquad
 \Path A'_\mu(L)\longrightarrow\Path A'_{\mu+1}(L)
\end{equation}
is an h-cofibration for all \(L\in\Obj(\Pcat)\).

The assertion at \(0\) is the hypothesis on \(s\).  At a successor stage there are three cases.
\begin{enumerate}
\item A pushout of \(C\) adds one isolated state and creates no path.  Thus
the maps in \eqref{eq:successor-hcof} are identities, and \(s_{\mu+1}\) is a
bijection on states with the same componentwise path maps as \(s_\mu\).
\item A pushout of \(R\) identifies two states and freely creates exactly the
reduced words of Lemma~\ref{lem:R-step}.  That lemma proves both the
induction assertion for \(s_{\mu+1}\) and the two h-cofibration assertions.
\item A pushout of a generating globe is covered by
Proposition~\ref{prop:globe-step}, which again proves both assertions.
\end{enumerate}

Let \(\mu\) be a limit ordinal and suppose that the induction is established below \(\mu\).  By \eqref{eq:successor-hcof} and \cite[Proposition~2.6]{leftproperflow}, both path towers consist objectwise of relative-\(T_1\) inclusions.  The path functor therefore commutes with their colimits:
\begin{equation}\label{eq:path-limit}
 \liminj_{\nu<\mu}\Path A_\nu\xrightarrow{\cong}\Path A_\mu,
 \qquad
 \liminj_{\nu<\mu}\Path A'_\nu\xrightarrow{\cong}\Path A'_\mu
\end{equation}
by \cite[Theorem~6.14]{Moore1}.  For every \(L\in\Obj(\Pcat)\), Corollary~\ref{cor:transfinite-gluing}, applied to the evaluated towers, shows that the map between the two colimits in \eqref{eq:path-limit} is a weak homotopy equivalence.  On states, the natural transformation below \(\mu\) is objectwise a bijection, hence a natural isomorphism of diagrams of sets; its colimit is again a bijection.  It follows that \(s_\mu\) is a weak equivalence.  This completes the induction.  In particular, the base change
\[
 \bar s=s_\lambda: X\longrightarrow
 A'\mathop{\coprod}_{A}X
\]
is a weak equivalence whenever \(i\) is a relative generating-cell complex.

Finally let \(i: A\to X\) be an arbitrary q-cofibration.  Apply the small object argument and the retract characterization of cofibrations \cite[Theorem~2.1.14 and Corollary~2.1.15]{MR99h:55031} to factor it as
\[
 A\xlongrightarrow{j}Y\xlongrightarrow{p}X,
\]
where \(j\) is a relative cell complex on \eqref{eq:Pflow-generators} and \(p\) is a trivial q-fibration.  Since \(i\) has the left lifting property with respect to \(p\), the square with top \(j\), left side \(i\), right side \(p\), and bottom \(\id_X\) has a lift \(r: X\to Y\).  Thus \(ri=j\) and \(pr=\id_X\), so \(i\) is a retract of \(j\) in the undercategory \(A\downarrow\PFlow\).  Pushout along \(s\) is a functor from this undercategory to \(A'\downarrow\PFlow\) and hence preserves that retract.  Consequently the induced pushout map for \(i\) is a retract, in the arrow category, of the induced pushout map for \(j\).  The latter is a weak equivalence by the cellular case.  A retract of a bijection is a bijection, and applying \(\pi_0\) and every based \(\pi_n\) shows that a retract of a weak homotopy equivalence is a weak homotopy equivalence.  Thus weak equivalences of \(\Pcat\)-flows are closed under retracts, so the pushout map for \(i\) is a weak equivalence.  This is exactly left properness.
\end{proof}

\section{Applications to interval reparametrization categories}
\label{sec:7}

The objects of \(\G\) and \(\M\) are the positive real numbers, with tensor product \(a\otimes b=a+b\).  Their enriched hom-spaces are
\begin{align*}
 \G(a,b)&=\{\text{nondecreasing homeomorphisms }[0,a]\to[0,b]\},\\
 \M(a,b)&=\{\text{nondecreasing continuous surjections }[0,a]\to[0,b]\}.
\end{align*}
They carry the \(\Delta\)-kelleyfication of the relative compact-open topology, and composition is composition of maps.  These data define reparametrization categories by \cite[Propositions~4.9 and~4.11]{Moore1}.  For \(\alpha\in\Pcat(a,a')\) and \(\beta\in\Pcat(b,b')\), where \(\Pcat\in\{\G,\M\}\), the \textit{block tensor} of maps is
\begin{equation}\label{eq:block}
 (\alpha\otimes\beta)(t)=
 \begin{cases}
  \alpha(t),&0\leq t\leq a,\\
  a'+\beta(t-a),&a\leq t\leq a+b.
 \end{cases}
\end{equation}
This is \cite[Definition~4.6]{Moore1}; its associativity and continuity
are part of \cite[Propositions~4.7, 4.9 and~4.11]{Moore1}.

We refer to \(\G\) and \(\M\) as the \textit{interval reparametrization categories}.

\begin{lem} \label{lem:inverse}
	The function 
	\[
	(-)^{-1}:\G(L,L') \to \G(L',L),\qquad f\longmapsto f^{-1}
	\]
	is continuous. 
\end{lem}

\begin{proof}
	This is \cite[Lemma~6.2]{Moore2}.
\end{proof}

\begin{thm}\label{thm:GM-cuttable}
The interval reparametrization categories \(\G\) and \(\M\) have cuts.
\end{thm}

\begin{proof}
Fix \(a,b,L>0\).  In both categories, choose
\begin{equation}\label{eq:GM-delta}
 \delta_L:[0,L]\longrightarrow[0,2L],\qquad
 \delta_L(t)=2t.
\end{equation}
This is a nondecreasing homeomorphism.  It therefore belongs to \(\G(L,2L)\), and, being surjective, to \(\M(L,2L)\) as well.

For \(\psi\in\Pcat(L,a+b)\), define the two cuts as follows.  If \(\Pcat=\G\), let \(\tau=\psi^{-1}(a)\) and put
\begin{align}
 c^{a,b}_1(\psi)(t)
  &=\psi\left(\frac{\tau t}{L}\right),                \label{eq:G-cut1}\\
 c^{a,b}_2(\psi)(t)
  &=\psi\left(\tau+\frac{(L-\tau)t}{L}\right)-a.     \label{eq:G-cut2}
\end{align}
The endpoint values and strict monotonicity show that these maps belong to \(\G(L,a)\) and \(\G(L,b)\), respectively.  If \(\Pcat=\M\), put
\begin{align}
 c^{a,b}_1(\psi)(t)&=\min\{\psi(t),a\},                \label{eq:M-cut1}\\
 c^{a,b}_2(\psi)(t)&=\max\{\psi(t)-a,0\}.             \label{eq:M-cut2}
\end{align}
These maps are continuous and nondecreasing.  Their endpoint values are, respectively, \((0,a)\) and \((0,b)\), so the intermediate value theorem makes them surjective.  They therefore belong to \(\M(L,a)\) and \(\M(L,b)\).

The topological spaces \(\Pcat(L,a+b)\), \(\Pcat(L,a)\), and \(\Pcat(L,b)\) are sequential because they are \(\Delta\)-generated. By \cite[Proposition~3]{Moore3}, convergence of sequences in each of these mapping spaces is equivalent to pointwise convergence. The formulas defining the four cut operators, together with Lemma~\ref{lem:inverse} for \(\G\) and the continuity of \(\min\) and \(\max\) for \(\M\), show that these operators preserve convergent sequences. They are therefore sequentially continuous and, since their domains are sequential, continuous.


We next verify block equivariance.  For \(\G\), if \(\chi=(\alpha\otimes\beta)\psi\), then the inverse image of \(a'\) under \(\alpha\otimes\beta\) is the singleton \(\{a\}\).  Therefore \(\chi^{-1}(a')=\psi^{-1}(a)\), and restricting \eqref{eq:block} to the two sides of this split point gives
\begin{align}
 c^{a',b'}_1((\alpha\otimes\beta)\psi)
   &=\alpha c^{a,b}_1(\psi),                          \label{eq:GM-equiv1}\\
 c^{a',b'}_2((\alpha\otimes\beta)\psi)
   &=\beta c^{a,b}_2(\psi).                          \label{eq:GM-equiv2}
\end{align}
For \(\M\), define
\[
 q^{a,b}_1(t)=\min\{t,a\},\qquad
 q^{a,b}_2(t)=\max\{t-a,0\}.
\]
The cuts are \(q_i^{a,b}\psi\).  A direct check on \([0,a]\) and \([a,a+b]\), using \eqref{eq:block} and the endpoint values of
\(\alpha,\beta\), gives
\[
 q^{a',b'}_1(\alpha\otimes\beta)=\alpha q^{a,b}_1,
 \qquad
 q^{a',b'}_2(\alpha\otimes\beta)=\beta q^{a,b}_2.
\]
Composition with \(\psi\) proves \eqref{eq:GM-equiv1}--\eqref{eq:GM-equiv2} for \(\M\) as well.

It remains to construct the normalization homotopy.  Let
\begin{equation}\label{eq:GM-N}
 N^{a,b}_L(\psi)=
 \bigl(c^{a,b}_1(\psi)\otimes c^{a,b}_2(\psi)\bigr)\delta_L.
\end{equation}
It is continuous by continuity of the cuts, composition, and the tensor of maps.  Equations \eqref{eq:GM-equiv1}--\eqref{eq:GM-equiv2} and functoriality of \(\otimes\) give, for \(k=\alpha\otimes\beta\),
\begin{equation}\label{eq:GM-N-equiv}
 N^{a',b'}_L(k\psi)=kN^{a,b}_L(\psi).
\end{equation}
For \(u\in[0,1]\), put
\begin{align}
 H_u(\psi)(t)
  &=\max\{\psi(t),N^{a,b}_L(\psi)(ut)\},             \label{eq:GM-H}\\
 K_u(\psi)(t)
  &=\max\{N^{a,b}_L(\psi)(t),\psi(ut)\}.            \label{eq:GM-K}
\end{align}
For \(\M\), these maps are continuous and nondecreasing, and their endpoint values are \(0,a+b\); the intermediate value theorem makes them surjective. Thus they belong to \(\M(L,a+b)\).

For \(\G\), they are homeomorphisms.  If \(u>0\), both functions entering either maximum are strictly increasing.  The maximum of two strictly increasing real functions is strictly increasing: if \(s<t\), whichever function realizes the maximum at \(s\) has a strictly larger value at \(t\), which is at most the maximum at \(t\).  If \(u=0\), then \(H_0(\psi)=\psi\) and \(K_0(\psi)=N^{a,b}_L(\psi)\).  In every case their endpoint values are \(0,a+b\), so they are homeomorphisms of intervals and belong to \(\G(L,a+b)\).

The two families are continuous.  Their adjoints are continuous by \eqref{eq:GM-H}--\eqref{eq:GM-K}, continuity of evaluation, of \((u,t)\mapsto ut\), and of the maximum map \(\mathbb R^2\to\mathbb R\). Cartesian closedness of \(\Top\) gives continuous maps into the internal mapping space.  Since the source is \(\Delta\)-generated and the images lie in the indicated relative subspaces, the universal property of \(\Delta\)-kelleyfication gives continuous maps into \(\Pcat(L,a+b)\); see \cite[Section~2]{Moore1}.

Every nondecreasing map \(k\) between intervals satisfies
\begin{equation}\label{eq:max-equiv}
 k(\max\{v,w\})=\max\{k(v),k(w)\}.
\end{equation}
Equations \eqref{eq:GM-N-equiv} and \eqref{eq:max-equiv} imply
\begin{equation}\label{eq:GM-HK-equiv}
 H_u(k\psi)=kH_u(\psi),\qquad K_u(k\psi)=kK_u(\psi).
\end{equation}
Moreover,
\[
 H_0(\psi)=\psi,\qquad
 H_1(\psi)=K_1(\psi)=\max\{\psi,N^{a,b}_L(\psi)\},\qquad
 K_0(\psi)=N^{a,b}_L(\psi).
\]
Define
\begin{equation}\label{eq:GM-Theta}
 \Theta^{a,b}_L(u,\psi)=
 \begin{cases}
  H_{2u}(\psi),&0\leq u\leq\frac12,\\
  K_{2-2u}(\psi),&\frac12\leq u\leq1.
 \end{cases}
\end{equation}
The two formulas agree at \(u=1/2\), so \(\Theta^{a,b}_L\) is continuous by the finite closed-cover pasting lemma.  For completeness, that lemma follows because the inverse image of a closed set under the pasted map is the union of its two inverse images, each closed in one member of the closed cover and hence closed in the whole source.  The displayed endpoint formulas show that \(\Theta^{a,b}_L\) joins the identity to \(N^{a,b}_L\), and \eqref{eq:GM-HK-equiv} gives its block equivariance.  Thus all the axioms of Definition~\ref{def:cuttable} hold for both \(\G\) and \(\M\).
\end{proof}

\begin{cor}\label{cor:GM-left-proper}
The q-model structures of \(\G\)-flows and \(\M\)-flows are left proper.
\end{cor}

\begin{proof}
Combine Theorem~\ref{thm:GM-cuttable} with Theorem~\ref{thm:left-proper}.
\end{proof}

\section{The case of the final category}
\label{sec:8}

Let \(\One\) be the category with one object \(*\) and one morphism \(\id_*\).  Its enriched hom-space is the one-point space, and its unique strict semimonoidal structure is given by \(*\otimes*=*\) and \(\id_*\otimes\id_*=\id_*\).  It is final because every \(\Top\)-enriched category admits a unique enriched functor to \(\One\).

\begin{proposition}\label{prop:one-cuttable}
The final category \(\One\) is a reparametrization category with cuts.
\end{proposition}

\begin{proof}
The semimonoidal structure is strict, the unique hom-space is contractible, and the tensor-factorization axiom has the unique solution \(\id_*=\id_*\otimes\id_*\).  Hence \(\One\) is a reparametrization category. There is only one possible choice for \(\delta_*\) and for each cut operator; all are \(\id_*\).  The normalization \(N^{*,*}_*\) is the identity of the one-point hom-space, and the constant homotopy is block-equivariant.  These data satisfy Definition~\ref{def:cuttable}.
\end{proof}

An enriched presheaf on \(\One\) is just a space.  Formula \eqref{eq:coend} has one summand and only identity relations, so it gives the natural homeomorphism
\[
 (D\otimes E)(*)\cong D(*)\times E(*).
\]
Consequently, a \(\One\)-flow is precisely a small topologically enriched semicategory, i.e., an ordinary flow.  This identification is also \cite[Definition~17]{Moore3}.

\begin{cor} \label{cor:one-left-proper}
The q-model structure of ordinary flows is left proper.
\end{cor}

\begin{proof}
Under the preceding identification, apply Proposition~\ref{prop:one-cuttable} and Theorem~\ref{thm:left-proper}.  This recovers \cite[Theorem~5.6]{leftproperflow}.
\end{proof}



\section{A symmetric strictly unital non-final example}
\label{sec:9}



Here \emph{unital} refers to the semimonoidal tensor product, not to the identity morphisms of the underlying category. Since the tensor product in a reparametrization category is strictly associative, a unit for a \(\Top\)-enriched semimonoidal category \(\mathcal C\) would consist of an object \(e\) and enriched natural isomorphisms
\[
\lambda_X: e\otimes X\xrightarrow{\cong}X,
\qquad
\rho_X: X\otimes e\xrightarrow{\cong}X
\]
satisfying the monoidal unit coherence; in particular, \(\id_X\otimes\lambda_Y=\rho_X\otimes\id_Y\) after using strict associativity to identify the two sources.

We first establish that the interval reparametrization categories \(\G\) and \(\M\) are nonunital. The absence of a strict unit is already visible on objects: a strict unit \(e>0\) would have to satisfy \(e+L=L\).  This observation alone would not exclude a non-strict monoidal unit, because any two objects \(a,b>0\) are isomorphic in both categories.  Indeed, the linear homeomorphism \(t\mapsto bt/a\) belongs to \(\G(a,b)\), its inverse is the corresponding linear homeomorphism from \([0,b]\) to \([0,a]\), and both maps also belong to \(\M\).

\begin{proposition}\label{prop:GM-nonunital}
	Neither \(\G\) nor \(\M\) is unital.  More strongly, for \(\Pcat\in\{\G,\M\}\) there do not exist an object \(e\) and an ordinary natural isomorphism
	\[
	\lambda_L: e\otimes L\xrightarrow{\cong}L
	\qquad(L\in\Obj(\Pcat)).
	\]
\end{proposition}

\begin{proof}
	Suppose that such an object \(e\) and such a natural isomorphism \(\lambda\) exist.  By definition of \(\G\) and \(\M\), \(e\) is a positive real number and \(e\otimes L=e+L\) for every \(L>0\).
	
	Fix \(L>0\).  The component \(\lambda_L\) is represented by a strictly increasing homeomorphism
	\[
	\lambda_L:[0,e+L]\longrightarrow[0,L].
	\]
	For \(\Pcat=\G\), this follows directly from the definition of its hom-spaces.  For \(\Pcat=\M\), the inverse of the isomorphism \(\lambda_L\) is a morphism whose composite with \(\lambda_L\) in both orders is an identity; hence the underlying continuous map of \(\lambda_L\) is bijective.  A continuous bijection from a compact interval to a Hausdorff interval is a homeomorphism: it is closed because the image of every closed, hence compact, subset is compact and therefore closed.  Finally, an injective nondecreasing map of intervals is strictly increasing.  This proves the claim in both cases.
	
	Put
	\[
	a=\lambda_L(e).
	\]
	Since \(0<e<e+L\) and \(\lambda_L\) is strictly increasing with endpoint values \(\lambda_L(0)=0\) and \(\lambda_L(e+L)=L\), we have
	\begin{equation}\label{eq:unitor-interior-point}
		0<a<L.
	\end{equation}
	Consider the endomorphism
	\[
	\phi_L:[0,L]\longrightarrow[0,L],
	\qquad
	\phi_L(t)=\frac{t^2}{L}.
	\]
	It is continuous and strictly increasing, has endpoint values \(0\) and \(L\), and has the continuous inverse \(t\mapsto\sqrt{Lt}\).  Thus \(\phi_L\in\G(L,L)\); being also a continuous nondecreasing surjection, it belongs to \(\M(L,L)\) as well.
	
	Naturality of \(\lambda\) with respect to \(\phi_L\) gives the commutative square
	\[
	\begin{tikzcd}
		e+L \arrow[r,"\id_e\otimes\phi_L"]
		\arrow[d,"\lambda_L"']
		& e+L \arrow[d,"\lambda_L"] \\
		L \arrow[r,"\phi_L"'] & L .
	\end{tikzcd}
	\]
	By the block formula \eqref{eq:block}, \((\id_e\otimes\phi_L)(e)=e\).  Evaluation of the commutative square at \(e\) therefore yields
	\[
	\phi_L(a)=\phi_L(\lambda_L(e))
	=\lambda_L((\id_e\otimes\phi_L)(e))
	=\lambda_L(e)=a.
	\]
	On the other hand, \eqref{eq:unitor-interior-point} gives
	\[
	\phi_L(a)=\frac{a^2}{L}<a,
	\]
	a contradiction.  Hence no natural left unitor exists.  A fortiori, neither \(\G\) nor \(\M\) admits a monoidal unit.
\end{proof}

The final category is an example of a symmetric strictly unital reparametrization category with cuts. We now exhibit an example of a symmetric strictly unital reparametrization category with cuts which is not equivalent to the final category. Let \(\Pmult\) have one object \(*\) and enriched endomorphism space
\[
\Pmult(*,*)=[0,1].
\]
The identity is \(1\), and composition is ordinary multiplication:
\[
t\circ s=ts \qquad (s,t\in[0,1]).
\]
This is a small \(\Top\)-enriched category.  Indeed, multiplication is a continuous map \([0,1]^2\to[0,1]\), it is associative, and \(1\) is its two-sided identity.

\begin{thm}\label{thm:example}
	There is a strict symmetric monoidal structure and a cut structure on \(\Pmult\) for which \(\Pmult\) is a symmetric reparametrization category with cuts.  Moreover, \(\Pmult\) is not equivalent to \(\One\).
\end{thm}

\begin{proof}
	There is only one possible tensor product on objects, namely \(*\otimes *=*\).  On morphisms define
	\begin{equation}\label{eq:tensor-multiplication}
		s\otimes t=st.
	\end{equation}
	It is continuous.  It preserves the identity since \(1\otimes1=1\), and it is compatible with composition: for \(s_1,s_2,t_1,t_2\in[0,1]\),
	\[
	(s_2s_1)\otimes(t_2t_1)
	=s_2s_1t_2t_1
	=s_2t_2s_1t_1
	=(s_2\otimes t_2)(s_1\otimes t_1).
	\]
	The middle equality uses commutativity of multiplication.  Thus \(\otimes\) is an enriched bifunctor.  Its strict associativity follows from \((rs)t=r(st)\), so \(\Pmult\) is a strict enriched semimonoidal category.
	
	The unique object \(\ast\) is a strict tensor unit.  Indeed, for every
	\(s\in[0,1]\),
	\[
	1\otimes s=s=s\otimes1.
	\]
	Hence both unitors are identity enriched natural transformations, and all unit coherence identities hold strictly.
	
	The only hom-space \([0,1]\) is contractible: the map
	\[
	H:[0,1]\times[0,1]\longrightarrow[0,1],
	\qquad H(u,s)=(1-u)s+u
	\]
	is a homotopy from the identity to the constant map with value \(1\). It remains to check the factorization axiom.  Both the source and the target have only the decomposition \(*=*\otimes *\).  For any
	\(\phi\in\Pmult(*,*)\), take \(\phi_1=\phi\) and \(\phi_2=1\).  Then
	\[
	\phi_1\otimes\phi_2=\phi\cdot1=\phi.
	\]
	Hence \(\Pmult\) is a reparametrization category.
	
	Define the symmetry by \(\gamma_{*,*}=1\).  For \(s,t\in[0,1]\), its naturality condition is
	\[
	\gamma_{*,*}(s\otimes t)=(t\otimes s)\gamma_{*,*},
	\]
	which reduces to \(st=ts\). The two hexagon identities expressing compatibility of the braiding with the associator \cite[Definition~8.1.1]{Tensor_categories}, as well as the symmetry identity \(\gamma_{\ast,\ast}^2=\id_{\ast}\)\cite[Definition~8.1.12]{Tensor_categories}, reduce to \(1=1\).  Therefore this symmetry satisfies all the axioms recalled above.
	
	We construct a cut structure.  Put
	\[
	\delta_*=1,
	\qquad c^{*,*}_1(\psi)=0,
	\qquad c^{*,*}_2(\psi)=0
	\quad(\psi\in[0,1]).
	\]
	The two cut operators are continuous because they are constant.  If \(\alpha,\beta,\psi\in[0,1]\), then, using \eqref{eq:tensor-multiplication} and composition by multiplication,
	\[
	c^{*,*}_1((\alpha\otimes\beta)\psi)=0
	=\alpha 0=\alpha c^{*,*}_1(\psi),
	\qquad
	c^{*,*}_2((\alpha\otimes\beta)\psi)=0
	=\beta 0=\beta c^{*,*}_2(\psi).
	\]
	Thus the block-equivariance identities \eqref{eq:cut-equiv1}--\eqref{eq:cut-equiv2} hold.  The normalization is
	\[
	N^{*,*}_*(\psi)=(0\otimes0)\delta_*=0.
	\]
	Define
	\begin{equation}\label{eq:theta-example}
		\Theta^{*,*}_*(u,\psi)=(1-u)\psi.
	\end{equation}
	This map is continuous, and its values at \(u=0\) and \(u=1\) are respectively \(\psi\) and \(0=N^{*,*}_*(\psi)\).  Finally,
	\begin{align*}
		\Theta^{*,*}_*(u,(\alpha\otimes\beta)\psi)
		&=(1-u)\alpha\beta\psi \\
		&=\alpha\beta(1-u)\psi \\
		&=(\alpha\otimes\beta)\Theta^{*,*}_*(u,\psi),
	\end{align*}
	so \eqref{eq:Theta-equiv} holds.  This proves that \(\Pmult\) has cuts.
	
	It remains to prove the asserted non-equivalence.  Any equivalence of ordinary categories is faithful: if \(F:\mathcal C\to\mathcal D\) has a quasi-inverse \(G\) and a natural isomorphism \(\eta:\id_{\mathcal C}\Rightarrow GF\), then \(F(f)=F(g)\) implies
	\[
	f=\eta_y^{-1}GF(f)\eta_x
	=\eta_y^{-1}GF(g)\eta_x=g
	\]
	for all parallel arrows \(f,g: x\to y\).  The underlying ordinary category of \(\Pmult\) has the distinct endomorphisms \(0\) and \(1\), whereas the underlying category of \(\One\) has only one endomorphism.  The unique functor \(\Pmult\to\One\) therefore sends \(0\) and \(1\) to the same arrow and is not faithful; consequently it is not an equivalence.  An enriched or a symmetric semimonoidal equivalence would in particular induce an equivalence of the underlying ordinary categories.  Hence no such equivalence between \(\Pmult\) and \(\One\) exists.
\end{proof}



\begin{thebibliography}{10}
	
	\bibitem{Barthel-Riel}
	T.~Barthel and E.~Riehl.
	\newblock On the construction of functorial factorizations for model
	categories.
	\newblock {\em Algebr. Geom. Topol.}, 13(2):1089--1124, 2013.
	\newblock \href {https://doi.org/10.2140/agt.2013.13.1089}
	{\path{https://doi.org/10.2140/agt.2013.13.1089}}.
	
	\bibitem{hocolimfacile}
	D.~Dugger and D.~C. Isaksen.
	\newblock Topological hypercovers and {\(\mathbb A^1\)}-realizations.
	\newblock {\em Math. Z.}, 246(4):667--689, 2004.
	\newblock \href {https://doi.org/10.1007/s00209-003-0607-y}
	{\path{https://doi.org/10.1007/s00209-003-0607-y}}.
	
	\bibitem{Tensor_categories}
	P.~Etingof, S.~Gelaki, D.~Nikshych, and V.~Ostrik.
	\newblock {\em Tensor categories}, volume 205 of {\em Math. Surv. Monogr.}
	\newblock Providence, RI: American Mathematical Society (AMS), 2015.
	\newblock \href {https://doi.org/10.1090/surv/205}
	{\path{https://doi.org/10.1090/surv/205}}.
	
	\bibitem{DAT_book}
	L.~Fajstrup, E.~Goubault, E.~Haucourt, S.~Mimram, and M.~Raussen.
	\newblock {\em Directed algebraic topology and concurrency. {With} a foreword
		by {Maurice} {Herlihy} and a preface by {Samuel} {Mimram}}.
	\newblock SpringerBriefs Appl. Sci. Technol. Springer, 2016.
	\newblock \href {https://doi.org/10.1007/978-3-319-15398-8}
	{\path{https://doi.org/10.1007/978-3-319-15398-8}}.
	
	\bibitem{FajstrupRaussenGoubault}
	L.~Fajstrup, M.~Raussen, and E.~Goubault.
	\newblock Algebraic topology and concurrency.
	\newblock {\em Theor. Comput. Sci.}, 357(1-3):241--278, 2006.
	\newblock \href {https://doi.org/10.1016/j.tcs.2006.03.022}
	{\path{https://doi.org/10.1016/j.tcs.2006.03.022}}.
	
	\bibitem{dgrtop}
	P.~Gaucher.
	\newblock Enriched diagrams of topological spaces over locally contractible
	enriched categories.
	\newblock {\em New York J. Math.}, 25:1485--1510, 2019.
	
	\bibitem{Moore1}
	P.~Gaucher.
	\newblock Homotopy theory of {M}oore flows ({I}).
	\newblock {\em {Compositionality}}, 3(3), 2021.
	\newblock \href {https://doi.org/10.32408/compositionality-3-3}
	{\path{https://doi.org/10.32408/compositionality-3-3}}.
	
	\bibitem{Moore2}
	P.~Gaucher.
	\newblock Homotopy theory of {M}oore flows ({II}).
	\newblock {\em {Extr. Math.}}, 36(2):157--239, 2021.
	\newblock \href {https://doi.org/10.17398/2605-5686.36.2.157}
	{\path{https://doi.org/10.17398/2605-5686.36.2.157}}.
	
	\bibitem{leftproperflow}
	P.~Gaucher.
	\newblock Left properness of flows.
	\newblock {\em Theory Appl. Categ.}, 37(19):562--612, 2021.
	
	\bibitem{Moore3}
	P.~Gaucher.
	\newblock Homotopy theory of {M}oore flows ({III}).
	\newblock {\em North-West. Eur. J. Math.}, 10:55--113, 2024.
	
	\bibitem{GoubaultUserGuide}
	E.~Goubault.
	\newblock Geometry and concurrency: {A} user's guide.
	\newblock {\em Math. Struct. Comput. Sci.}, 10(4):411--425, 2000.
	\newblock \href {https://doi.org/10.1017/S0960129500003133}
	{\path{https://doi.org/10.1017/S0960129500003133}}.
	
	\bibitem{Grandis}
	M.~Grandis.
	\newblock {\em Directed algebraic topology. {Models} of non-reversible worlds},
	volume~13 of {\em New Math. Monogr.}
	\newblock Cambridge: Cambridge University Press, 2009.
	\newblock \href {https://doi.org/10.1017/CBO9780511657474}
	{\path{https://doi.org/10.1017/CBO9780511657474}}.
	
	\bibitem{ref_model2}
	P.~S. Hirschhorn.
	\newblock {\em Model categories and their localizations}, volume~99 of {\em
		Mathematical Surveys and Monographs}.
	\newblock American Mathematical Society, Providence, RI, 2003.
	\newblock \href {https://doi.org/10.1090/surv/099}
	{\path{https://doi.org/10.1090/surv/099}}.
	
	\bibitem{MR99h:55031}
	M.~Hovey.
	\newblock {\em Model categories}.
	\newblock American Mathematical Society, Providence, RI, 1999.
	\newblock \href {https://doi.org/10.1090/surv/063}
	{\path{https://doi.org/10.1090/surv/063}}.
	
	\bibitem{Riehl}
	E.~Riehl.
	\newblock {\em Categorical homotopy theory}, volume~24 of {\em New Math.
		Monogr.}
	\newblock Cambridge: Cambridge University Press, 2014.
	\newblock \href {https://doi.org/10.1017/CBO9781107261457}
	{\path{https://doi.org/10.1017/CBO9781107261457}}.
	
	\bibitem{SchwanzlVogt}
	R.~Schw{\"a}nzl and R.~M. Vogt.
	\newblock Strong cofibrations and fibrations in enriched categories.
	\newblock {\em Arch. Math.}, 79(6):449--462, 2002.
	\newblock \href {https://doi.org/10.1007/PL00012472}
	{\path{https://doi.org/10.1007/PL00012472}}.
	
\end{thebibliography}


%\bibliographystyle{../plainurlwithoutprefixDOI} 
%\bibliography{../Bibliotheque} 



\end{document}
